\textbf{COLUNI 2025 -- Questão 04}\par\medskip Os excertos, a seguir, fazem referência a um elemento importante do enredo de \textit{Torto arado}: os encantados.

\par

O curador Zeca Chapéu Grande tudo podia. Se transformava em muitos encantados nas noites de jarê. Mudava a voz, cantava, rodopiava ágil pela sala, investido dos poderes dos espíritos das matas, das águas, das serras e do ar. Meu pai curava loucos e bêbados, colocaria meu pedaço de língua em minha boca. (VIEIRA JUNIOR, Itamar. \textit{Torto arado}. São Paulo: Todavia, 2019. p. 126)

\par

Meu cavalo morreu e não tenho mais montaria para caminhar como devo, da forma que um encantado deve se apresentar entre os homens, como deve aparecer por esse mundo. Desde então, passei a vagar sem rumo, arrodeando aqui, arrodeando acolá, procurando um corpo que pudesse me acolher. Meu cavalo era uma mulher chamada Miúda, mas quando me apossava de sua carne seu nome era Santa Rita Pescadeira. Foi nela que cavalguei por tempo, não conto o tempo, mas montei o corpo de Miúda, solitária. Sou muito mais antiga que os cem anos de Miúda. (VIEIRA JUNIOR, Itamar. \textit{Torto arado}. São Paulo: Todavia, 2019. p. 203)

\par

Sobre os \underline{encantados} constituintes da obra, é CORRETO afirmar que são:\begin{enumerate}
\item entidades ancestrais, ligadas à natureza, servindo de ponte entre o mundo material e espiritual.
\item figuras de uma tradição europeia, chegando ao Brasil com os colonizadores portugueses.
\item entidades da igreja católica, confirmando a constituição da religiosidade na obra.
\item figuras de importância ínfima para a relação entre os mundos espirituais e humanos.\end{enumerate}
